\chapter{Defining the Budget Relay}
\label{ch:budget-relay}

\begin{quotation}
\noindent\textit{Synopsis.} This chapter introduces the budget relay, the
monograph's primary new theoretical object: a chain of renderings connected
by independent local repair economies, each governed by its own costs,
priorities, and constraints. The chapter formalizes how distinctions are
discarded under local budget pressure and why such losses become
irreversible once no record of them remains. Root objects and
budget-aware loss ledgers are introduced as the mechanisms for measuring,
and in some cases recovering from, downstream insufficiency; preserving a
root object is shown to be fundamentally different from preserving every
downstream rendering, since it alone permits re-derivation under newly
discovered tasks.
\end{quotation}

This chapter carries disproportionate weight in the monograph relative to
its length. The object defined in \cref{def:budget-relay} is the only
genuinely new mathematical object this book introduces; every later
theorem is, in an important sense, a property of this definition rather
than an independent contribution. If this definition is not judged
inevitable and useful on its own terms, no proof in
Chapter~\ref{ch:composition-failure} will rescue it. The chapter is
therefore written at a slower pace than its material would otherwise
require, and it deliberately stress-tests the definition against a fully
worked numerical instance before the rest of the monograph is permitted to
build on it. Two genuine gaps surface in the course of this stress-testing
--- the underdetermination of the local discard rule, and the bootstrapping
cost of preserving a root object --- and both are carried forward
explicitly into Appendix~\ref{app:open-questions} rather than quietly
absorbed into the definition.

\section{Why a Chain, and Not Just a Rendering}

Chapter~\ref{ch:distinctions} defined a rendering as a single projection
from a source's distinction space to a smaller one, and
Chapter~\ref{ch:admissibility} asked, of a single such rendering, whether
it is sufficient for a given task. Neither chapter had reason to ask what
happens when a rendering is itself produced from an earlier rendering, and
that one from a still earlier one, with a different party responsible for
each step and no party responsible for the sequence as a whole. This is
not a generalization added for its own sake. It is the ordinary condition
of almost every substrate a real consumer actually encounters: a
transcript is produced from a recording, not from the recorded event
itself; a textbook's account of a proof is produced from the proof, not
independently derived; a legal memorandum summarizing a contract is
produced from the contract, which was itself drafted from negotiation
notes that no longer exist. Treating each of these as an isolated
rendering of some single, fixed original -- as Chapters~\ref{ch:distinctions}
and~\ref{ch:admissibility} implicitly do when they speak of ``the''
discard set of ``a'' rendering $r$ of ``an'' object $e$ -- silently assumes
that $e$ itself, and not some intermediate rendering of it, is always
available to compare against. That assumption is false for essentially
every substrate chain of practical interest, and its falsity is precisely
what Chapter~\ref{ch:synchronic-to-diachronic} identified as the reason
the repair-budget apparatus of Chapter~\ref{ch:repair-economics} does not,
by itself, resolve the continuation problem of
Chapter~\ref{ch:continuation-problem}.

\section{The Substrate Chain}

\begin{definition}[Substrate Chain]
\label{def:substrate-chain}
A substrate chain is a sequence
\[
A_0 \xrightarrow{B_0} A_1 \xrightarrow{B_1} \cdots \xrightarrow{B_{n-1}} A_n
\]
where each $A_i$ is a rendering (\cref{def:rendering}) and each transition
is governed by a local budget triple $B_i = (C_i, I_i, \beta_i)$: a local
repair cost function, a local interference function, and the budget
available at that stage.
\end{definition}

Two restrictions are built into \cref{def:substrate-chain} and are worth
stating openly rather than leaving implicit. First, the chain is linear
and finite: each $A_i$ has exactly one predecessor and one successor, and
there are finitely many stages. The branching generalization, in which a
single $A_i$ gives rise to several independent successors, is deferred to
Chapter~\ref{ch:fanout}, where it is shown to introduce a genuinely new
failure mode (silent cross-branch disagreement) rather than being a
straightforward extension of what follows here. Whether a substrate chain
with infinitely many stages behaves differently in kind, and not merely in
degree, from the finite case treated throughout this monograph is left
unaddressed and is recorded as an open question in
Appendix~\ref{app:open-questions}. Second, $A_0$ is not assumed to be the
original object $e$ of Chapter~\ref{ch:distinctions}; it is only assumed to
be the earliest rendering the chain's analysis takes as given. A substrate
chain is therefore always relative to a choice of starting point, and
nothing in this chapter's machinery requires that starting point to be
loss-free relative to whatever preceded it.

\subsection{Local Budgets as Stage-Local Repair Economies}

The triple $B_i = (C_i, I_i, \beta_i)$ is not an unrelated new primitive.
It is Chapter~\ref{ch:repair-economics}'s repair-economic apparatus,
restricted to a single stage and to that stage's own repair-stable domain.
Concretely, $C_i$ is the restriction of the repair cost function
$\repcost{\cdot}$ of \cref{def:sharpness} to the distinctions present at
stage $i$; $I_i$ is the restriction of the interference functional
$\interf{\cdot}{\cdot}$ of \cref{def:interference} to pairs of distinctions
competing for stage $i$'s budget; and $\beta_i$ plays the role, at stage
$i$ alone, that the total budget $\Phi_{\mathrm{tot}}$ plays for a single
observer in \cref{thm:redistribution}. The relationship between $\beta_i$
and $\Phi_{\mathrm{tot}}$ is best understood by contrast rather than by
analogy: $\Phi_{\mathrm{tot}}$ is conserved across an observer's entire,
simultaneous distinction space, so that \cref{thm:redistribution} can speak
of resources being redistributed \emph{within} it. No such global quantity
is conserved across a substrate chain's stages. $\beta_0, \beta_1, \ldots,
\beta_{n-1}$ are, in general, entirely independent numbers set by
independent parties for independent reasons, and nothing in
\cref{def:substrate-chain} relates $\beta_i$ to $\beta_j$ for $i \neq j$.
This is the formal expression of the distinction drawn informally in
Chapter~\ref{ch:synchronic-to-diachronic}: the Redistribution Theorem
governs one $\Phi_{\mathrm{tot}}$ shared by many distinctions; a substrate
chain has many $\beta_i$, each shared by nothing beyond its own stage.

\section{Budget-Constrained Discard, and Its Underdetermination}

\begin{definition}[Budget-Constrained Discard]
\label{def:budget-discard}
A distinction $d$ present in $A_i$ is dropped in $A_{i+1}$ if maintaining
$d$ exceeds $\beta_i$, or if maintaining $d$ interferes too strongly, per
$I_i$, with a higher-priority distinction competing for the same stage's
budget.
\end{definition}

\cref{def:budget-discard} is deliberately schematic, and it is worth being
honest, before building anything further on it, about exactly what the
schema leaves open. The phrase ``higher-priority distinction'' presupposes
a priority ordering over the distinctions present at stage $i$, and
\cref{def:budget-discard} does not say where that ordering comes from or
what happens when two admissible orderings would license different
discards under the same budget $\beta_i$. This is not a minor omission to
be patched in passing. Two natural instantiations of
\cref{def:budget-discard} exist, they are not equivalent, and the
monograph's later theorems need to be checked against both rather than
silently assumed to hold for whichever one is easiest to compute with.

\begin{construction}[Greedy Priority Discard]
\label{con:greedy-discard}
Fix a priority ordering $d_{(1)} \succ d_{(2)} \succ \cdots \succ d_{(m)}$
over the distinctions present at stage $i$, reflecting their importance to
the stage's own task $\tau_i$. Process the distinctions in this order,
maintaining a running total of committed budget. A distinction is funded,
and hence preserved into $A_{i+1}$, if the remaining budget after funding
every higher-priority distinction is at least $C_i(d)$; otherwise it is
dropped, and processing continues to the next distinction in the ordering
without returning to reconsider a dropped distinction later.
\end{construction}

\begin{construction}[Cost-Optimal Discard]
\label{con:optimal-discard}
Given the same priority ordering interpreted now as a value assignment
$v(d_{(k)})$ decreasing in $k$, select the subset $S$ of stage $i$'s
distinctions maximizing $\sum_{d \in S} v(d)$ subject to $\sum_{d \in S}
C_i(d) \le \beta_i$, breaking ties in favor of higher-priority
distinctions.
\end{construction}

\begin{remark}[The Two Constructions Diverge]
\label{rem:discard-rules-diverge}
\cref{con:greedy-discard} and \cref{con:optimal-discard} are the discard
rule and the knapsack-optimal packing of the same underlying priorities and
costs, respectively, and they can select different subsets for the same
$(C_i, I_i, \beta_i)$ and the same priority ordering. A single
low-priority, low-cost distinction that the greedy rule would skip in
order to fund several higher-priority items may be exactly the distinction
a global optimization would sacrifice one mid-priority item to include,
whenever doing so raises the total value funded. Neither construction is
privileged by \cref{def:budget-discard} as stated; the definition is a
constraint on the \emph{outcome} of a stage's discard decision, not a
specification of the \emph{procedure} that produces it.
\end{remark}

This divergence is not a defect to be resolved before the theory can
proceed. It is the correct level of generality for a schema meant to cover
discard decisions made by a codec's rate-control algorithm, a human
editor's judgment, a court's interpretation of relevance, and an
institution's retention policy, none of which can be assumed to implement
the same procedure. What it does mean, and what is recorded explicitly
here so that it cannot be forgotten by the time Appendix~\ref{app:proof} is
drafted, is that any theorem claiming to hold for \emph{every} budget relay
--- including the Composition Failure Theorem of
Chapter~\ref{ch:composition-failure} --- must hold under both
\cref{con:greedy-discard} and \cref{con:optimal-discard}, and ideally under
any construction satisfying only the qualitative constraint of
\cref{def:budget-discard}. This robustness requirement is added to
Appendix~\ref{app:open-questions} as a condition Appendix~\ref{app:proof}'s
eventual proof must be checked against, not assumed to satisfy by
construction.

\section{The Budget Relay, Defined}

\begin{definition}[Budget Relay]
\label{def:budget-relay}
A substrate chain (\cref{def:substrate-chain}) is a \emph{budget relay} if,
for every stage $i$, the transition $B_i$ is selected using only the
current rendering $A_i$ and the stage's own local task $\tau_i$ --- that
is, without reference to the loss ledgers $L_0, \ldots, L_{i-1}$ of any
earlier stage, without reference to the budget triple $B_j$ of any other
stage $j \neq i$, and without appeal to any interference term spanning
distinctions held at different stages.
\end{definition}

This restates, and sharpens, an earlier and vaguer formulation of the same
condition. The earlier formulation spoke of stages having no ``shared
interference term or shared discard record\ldots by default,'' which left
unclear exactly what a stage is and is not permitted to consult when
setting its own budget. \cref{def:budget-relay} as now stated answers that
question directly: a stage may see the rendering handed to it and may know
its own task, and nothing else about the chain's history or future. This
is doing the real work of the definition, and it is worth dwelling on why.
A substrate chain in which every stage consults a common ledger and a
common, chain-wide interference functional before setting its own budget
is, for the purposes of this monograph's apparatus, indistinguishable from
a single observer's economy distributed across time: its total behavior is
governed by \cref{thm:redistribution} rather than by anything new. That
case is not excluded from consideration --- Chapter~\ref{ch:coordination}
treats exactly this kind of shared-ledger access as one point on a
spectrum of partial-coordination mechanisms --- but it is not the default
case, and \cref{def:budget-relay} is written to isolate the default case
cleanly: a relay in which coordination, if it exists at all, must be added
back in deliberately and explicitly, rather than assumed from the outset.

\begin{remark}[Degenerate Cases]
\label{rem:degenerate-relay}
A substrate chain with $n = 1$ satisfies \cref{def:budget-relay} vacuously,
since there is only one stage and no other stage's budget or ledger to
consult; this case reduces to the single rendering already treated in
Chapters~\ref{ch:distinctions}--\ref{ch:admissibility} and adds nothing new.
The genuinely new content of \cref{def:budget-relay} appears only once
$n \ge 2$, and the sharpest illustrations of that content, developed in
\cref{sec:worked-relay} below, use $n = 3$.
\end{remark}

\section{The Root-Object Move}

\cref{prop:task-unavailability} showed that sufficiency cannot be
certified at the time a substrate is produced, because the relevant future
task space is generally unknown then. Within a budget relay this
observation sharpens into a design question rather than remaining a mere
impossibility result: since no stage can know which of its discards a
future, unanticipated task will need reversed, is there anything a relay's
earliest stage can do that does not require anticipating the future task
space at all?

\begin{definition}[Root Object]
\label{def:root-object}
A distinguished rendering $A_0$ in a substrate chain, treated as the
reference point against which all downstream loss is measured. Where
possible, ingestion-time policy should optimize for re-derivability of
downstream renderings from $A_0$, rather than for sufficiency relative to
an as-yet-unknown $\tau$.
\end{definition}

\begin{proposition}[Root Recovery]
\label{prop:root-recovery}
If the root object $A_0$ survives unaltered, any downstream rendering
$A_i$ can be regenerated under an alternative budget allocation $B_i'$.
\end{proposition}

\begin{proofsketch}
By \cref{def:substrate-chain}, every $A_i$ for $i>0$ is generated from
$A_0$ via some composition of transitions. If $A_0$ persists, the
transition $B_i$ can be replaced with $B_i'$ and $A_i'$ recomputed
directly from $A_0$; the loss incurred under $B_i$ is therefore reversible
whenever $A_0$ is available. Without a preserved $A_0$, distinctions
discarded at any stage are permanently unavailable to every later stage,
since \cref{def:budget-discard} provides no mechanism for a distinction to
reappear once dropped.
\end{proofsketch}

\cref{prop:root-recovery} is easy to overstate, and it is worth being
precise about what it does not claim. It does not claim that preserving
$A_0$ is free. Preserving a root object is itself a repair decision --
$A_0$ must be maintained against whatever perturbations threaten it
(storage degradation, format obsolescence, simple deletion under
retention-policy pressure) -- and that maintenance itself has a cost that
must be funded by \emph{some} budget, decided by \emph{someone}, before any
downstream task requiring it is known to exist.

\begin{remark}[The Root-Preservation Bootstrapping Problem]
\label{rem:root-bootstrap}
Preserving $A_0$ so that \cref{prop:root-recovery} remains available is
itself an instance of \cref{prop:task-unavailability}: the decision to fund
$A_0$'s preservation must be made before the future task space that would
justify the expense is known, by definition of that task space being open.
An institution cannot, even in principle, calculate the return on
archiving a root object against tasks that do not yet exist. This does not
undermine \cref{def:root-object}'s recommendation so much as clarify its
status: root preservation is not a solution to the sufficiency problem of
\cref{ch:admissibility}, but a wager that re-derivability will be worth
more, in expectation, than whatever the preservation itself costs. Whether
this wager can be given a more precise decision-theoretic treatment,
rather than being left at the level of a design recommendation, is left
open in Appendix~\ref{app:open-questions}.
\end{remark}

\section{Cumulative Discard and Its Monotonicity}

The claim that loss compounds along an uncoordinated chain has been used
informally since Chapter~\ref{ch:synchronic-to-diachronic}. Making it
precise requires first defining what ``cumulative'' discard means for a
chain of more than one transition, since \cref{def:discard-set} was stated
only for a single rendering.

\begin{definition}[Cumulative Discard Set]
\label{def:cumulative-discard}
For a substrate chain and indices $i \le j$, write $D_k := \discard{A_k
\to A_{k+1}}$ for the discard set of the single transition $B_k$. The
cumulative discard set between stages $i$ and $j$ is
\[
D(A_i \Rightarrow A_j) \;:=\; \bigcup_{k=i}^{j-1} D_k.
\]
\end{definition}

\begin{proposition}[Monotonicity of Cumulative Discard]
\label{prop:monotonicity}
For any substrate chain in which the root object is not available for
recovery (\cref{prop:root-recovery} does not apply), and for any $i \le j
\le l$,
\[
D(A_i \Rightarrow A_j) \;\subseteq\; D(A_i \Rightarrow A_l).
\]
\end{proposition}

\begin{proofsketch}
Immediate from \cref{def:cumulative-discard}: $D(A_i \Rightarrow A_l)$ is a
union over a strictly larger index range than $D(A_i \Rightarrow A_j)$,
and a union over a superset of index terms contains the union over the
subset. No further argument is required once cumulative discard is defined
as a union of per-transition discard sets, since \cref{def:budget-discard}
supplies no operation by which a distinction can be re-added to a later
$A_l$ once absent from an intermediate $A_k$ for $k < l$.
\end{proofsketch}

This result was labeled a Theorem in an earlier draft of this chapter, and
it is worth recording explicitly why that label has been withdrawn.
\cref{prop:monotonicity}'s inequality is, once
\cref{def:cumulative-discard} is stated correctly, an immediate consequence
of set union rather than a substantive derivation, and the monograph's own
epistemic-labeling discipline reserves the word ``Theorem'' for results
whose proof carries independent weight. The genuine content of the claim
that loss compounds along a chain is not the inclusion itself but the
prior, harder-won observation that \cref{def:budget-discard} provides no
mechanism for reversal absent \cref{prop:root-recovery} -- and that
observation is a modeling commitment built into
\cref{def:substrate-chain}, not a theorem about it.

\section{The Loss Ledger, Upgraded}

\begin{definition}[Budget-Aware Loss Ledger]
\label{def:loss-ledger}
The loss ledger $L_i$ attached to transition $B_i$ records not merely the
discard set $D_i = \discard{A_i \to A_{i+1}}$, but the pair $L_i = (D_i,
B_i)$: the discarded distinctions \emph{together with} the local budget
that produced the discard. This allows a downstream consumer to
distinguish a distinction dropped because it was irrelevant to that
stage's own purpose from one dropped because that stage's budget was
simply small -- a distinction with different implications for whether the
loss might be recoverable elsewhere in the relay.
\end{definition}

\cref{def:loss-ledger} is what separates a budget relay's provenance
machinery from the plain loss ledger of Chapter~\ref{ch:distinctions}: it
is a coordination device, not merely a passive record, and its role as
such is developed fully in Chapter~\ref{ch:coordination}. It is worth
noting immediately, however, a limitation that \cref{def:budget-relay}
itself imposes on $L_i$'s usefulness: because a budget relay's stage $i+1$
does not consult $L_j$ for $j < i$ when setting $B_{i+1}$, the mere
existence of a loss ledger at every stage does not, by itself, prevent the
composition failures analyzed in Chapter~\ref{ch:composition-failure}. A
loss ledger is a necessary condition for a downstream consumer, arriving
after the fact and inspecting the whole chain, to diagnose what happened;
it is not, on its own, a mechanism by which an intermediate stage avoids
causing the problem in the first place.

\section{A Fully Worked Instance}
\label{sec:worked-relay}

The definitions above are abstract enough that their consequences are easy
to state and easy to misjudge. This section fixes one small, fully
explicit numerical relay and works every definition above against it by
hand, so that later chapters -- particularly the proof of
Chapter~\ref{ch:composition-failure} and its constructive witness in
Chapter~\ref{ch:pipelines} -- have a concrete instance to generalize from
rather than only an abstract schema.

Consider a source rendering $A_0$ carrying five distinctions, $d_1,
\ldots, d_5$, with repair costs $C(d_1) = 1$, $C(d_2) = 2$, $C(d_3) = 3$,
$C(d_4) = 2$, and $C(d_5) = 4$, and suppose interference among them is
negligible, so that $I \approx 0$ throughout and the greedy construction
(\cref{con:greedy-discard}) and the optimal construction
(\cref{con:optimal-discard}) can be compared directly against the same
numbers.

\subsection{Stage \texorpdfstring{$0 \to 1$}{0 to 1}}

Suppose stage $0$'s task $\tau_0$ ranks the distinctions by priority
$d_1 \succ d_2 \succ d_4 \succ d_5 \succ d_3$, and that $\beta_0 = 4$.
Applying \cref{con:greedy-discard}: $d_1$ is funded (cost $1$, remaining
budget $3$); $d_2$ is funded (cost $2$, remaining budget $1$); $d_4$
cannot be funded (cost $2 > 1$ remaining) and is dropped; $d_5$ cannot be
funded (cost $4 > 1$ remaining) and is dropped; $d_3$ cannot be funded
(cost $3 > 1$ remaining) and is dropped. Thus $A_1$ retains $\{d_1, d_2\}$
and $D_0 = \{d_3, d_4, d_5\}$. Applying \cref{con:optimal-discard} to the
same priorities and budget yields the identical subset in this instance,
since $\{d_1, d_2\}$ is in fact the value-maximizing feasible set here;
\cref{rem:discard-rules-diverge}'s divergence is real in general but does
not arise for every choice of numbers, and this instance was chosen partly
to establish a clean baseline before a divergent case is considered
elsewhere.

The loss ledger at this stage is $L_0 = (\{d_3, d_4, d_5\}, B_0)$, with
$B_0$ recording $\beta_0 = 4$ against the priority ordering above -- so
that a downstream consumer inspecting $L_0$ can see that $d_3$, $d_4$, and
$d_5$ were dropped because $\tau_0$ ranked them low and $\beta_0$ was
small relative to their combined cost, not because they were judged
worthless in any absolute sense.

\subsection{Stage \texorpdfstring{$1 \to 2$}{1 to 2}}

Suppose stage $1$'s task $\tau_1$ requires only $\{d_1, d_2\}$, exactly
what survived into $A_1$, and that $\beta_1 = 3$ -- exactly sufficient to
fund both at cost $1 + 2 = 3$. Then $A_2$ retains $\{d_1, d_2\}$ and $D_1 =
\emptyset$. By \cref{def:cumulative-discard}, $D(A_0 \Rightarrow A_2) =
D_0 \cup D_1 = \{d_3, d_4, d_5\}$, unchanged from $D(A_0 \Rightarrow A_1)$
-- an explicit instance of \cref{prop:monotonicity}'s inclusion, here
realized as equality rather than strict containment, since stage $1$
happened to discard nothing further.

\subsection{A Task Neither Stage Anticipated}

Suppose a task $\tau^*$ is later discovered, requiring recovery of $d_3$
-- some distinction, present in $A_0$, that neither $\tau_0$ nor $\tau_1$
had reason to prioritize. Since $d_3 \notin A_1$ and $A_2$ is derived only
from $A_1$, $d_3$ is unrecoverable from $A_2$ by
\cref{def:budget-discard}'s irreversibility and confirmed unrecoverable in
general by \cref{prop:monotonicity}. Neither $B_0$ nor $B_1$ was
irrational relative to its own task: $\beta_0$ was spent optimally against
$\tau_0$'s stated priorities, and $\beta_1$ was spent optimally, indeed
exactly, against $\tau_1$'s. $A_2$ is nonetheless insufficient
(\cref{def:sufficiency}) for $\tau^*$. This four-paragraph construction is
already, in miniature, the mechanism that Chapter~\ref{ch:composition-failure}
formalizes as Theorem~\ref{thm:composition-failure} and that
Chapter~\ref{ch:pipelines} exhibits again at full scale against a real
media pipeline; nothing about the general theorem's proof will do more
than what has already happened here with five distinctions and two
stages.

\subsection{Root Recovery, Demonstrated}

Suppose, finally, that $A_0$ itself was preserved as a root object
(\cref{def:root-object}) rather than discarded once $A_1$ was produced.
Upon discovering $\tau^*$, a new transition $B_0'$ can be computed directly
from $A_0$ using a revised priority ordering that ranks $d_3$ above $d_4$,
yielding a new rendering $A_1'$ retaining $\{d_1, d_2, d_3\}$ and
satisfying $\tau^*$. This is \cref{prop:root-recovery} made concrete: the
insufficiency exhibited above was permanent for the original chain
$A_0 \to A_1 \to A_2$, but not permanent for the underlying situation,
precisely because $A_0$, and not merely $A_1$ or $A_2$, remained available
to be re-derived from.

\section{What This Definition Commits Us To}

It is worth closing this chapter with an explicit account of exactly what
has been assumed, so that later chapters relaxing any one of these
assumptions can be read as deliberate extensions rather than silent
changes of subject. A substrate chain, as defined here, is linear rather
than branching (relaxed in Chapter~\ref{ch:fanout}); finite rather than
unbounded (left open in Appendix~\ref{app:open-questions}); governed at
each stage by a discard rule whose outcome, but not procedure, is
constrained by \cref{def:budget-discard} (with the resulting
underdetermination flagged in \cref{rem:discard-rules-diverge} as a
robustness condition the proofs of Chapter~\ref{ch:composition-failure}
must satisfy); irreversible in its losses except through the specific,
costly mechanism of root recovery (\cref{prop:root-recovery}, itself
subject to the bootstrapping problem of \cref{rem:root-bootstrap}); and
uncoordinated by default, in the precise sense given in
\cref{def:budget-relay}, with coordination treated in
Chapter~\ref{ch:coordination} as something that must be added back in
deliberately rather than assumed. Every one of these five commitments is a
modeling choice, not a discovered fact about the world, and the worked
instance of \cref{sec:worked-relay} was constructed specifically to make
each commitment's consequences visible in a case small enough to check by
hand before trusting the general theory built on top of it.
